\section{Christ, Mary, and the Angels}\label{sec:christ-mary-angels}

``And again, when he bringeth in the first begotten into the world, he saith: And let all the angels of God adore him'' (Heb 1:6). The Letter to the Hebrews opens by measuring the eternal Son against the entire celestial court---not as the highest archangel among created spirits, but as Creator, Heir, and Lord before whom every angel falls down in worship. St.~Athanasius seized upon this verse against the Arians (\emph{Against the Arians}~I.61--62): angels say ``I am sent unto thee'' and ``The Lord has commanded,'' whereas the Son receives their adoration and declares, ``I am in the Father, and the Father in Me.'' St.~Ambrose drew the same line for the Holy Spirit as well (\emph{On the Faith}~V.18.225; \emph{On the Holy Spirit}~III.16.142): Thrones, Dominations, Cherubim, and Seraphim serve and adore the Son, while Scripture never bids the Spirit worship, because the Spirit stands not among ministers but within the one Godhead. Yet Catholic theology presses the mystery one step deeper: St.~Thomas Aquinas demonstrates that Christ even \emph{as man} is Head of all the angels.

\subsection{Christ as Head of the Angels (\ST{III, q.~8, a.~4})}

Building upon St.~Paul's declaration that Christ ``is the head of all Principality and Power'' (Col 2:10), Aquinas asks in the \emph{Tertia Pars} whether Christ in his sacred humanity exercises genuine headship over the angelic choirs (\ST{III, q.~8, a.~4}). His proof starts from the unity of the heavenly city established in \ST{I, q.~108, a.~1} and a.~8. Wherever a multitude is ordered to one common end, it constitutes a single mystical body; ``both men and angels are ordained to one end, which is the glory of the Divine fruition. Hence the mystical body of the Church consists not only of men but of angels.'' Within this single Church of heaven and earth, Christ holds the place of Head even according to his assumed human nature, because ``He is nearer God, and shares His gifts more fully, not only than man, but even than angels; and of His influence not only men but even angels partake'' (Eph 1:20--22; Matt 4:11).

Three classic difficulties test this claim, and Aquinas's replies illuminate the deepest architecture of Christology and angelology together. First, because the Letter to the Hebrews states that Christ did not take hold of angelic nature but of the seed of Abraham (Heb 2:16), an objector might deny the conformity of nature required between head and members. Aquinas replies that Christ's headship over human beings touches chiefly their intellectual souls, ``wherein men agree with angels in generic nature, though not in specific nature'' (\ST{III, q.~8, a.~4, ad~1}). Second, if the Church is the congregation of the faithful walking by faith rather than sight (2 Cor 5:6--7), how can blessed angels belong to her? On earth the Church is the company of wayfarers, Aquinas answers, ``but, in heaven, it is the congregation of comprehensors. Now Christ was not merely a wayfarer, but a comprehensor. And therefore He is the Head not merely of the faithful, but of comprehensors, as having grace and glory most fully'' (\ST{III, q.~8, a.~4, ad~2}). Third, reconciling St.~Augustine's rule (\emph{Tractates on John}~XIX, XXIII) that bodily nature cannot act directly on a spirit, Aquinas appeals to the hypostatic union: ``the humanity of Christ, by virtue of the spiritual nature, i.e.\ the Divine, can cause something not only in the spirits of men, but also in the spirits of angels, on account of its most close conjunction with God, i.e.\ by personal union'' (\ST{III, q.~8, a.~4, ad~3}).

In this article the mature Aquinas surpasses his own commentary on Colossians. In \emph{Super Colossenses}~2:10, he had described Christ as Head of the angels principally as their King and Lord, reserving conformity of nature to human beings; in \ST{III, q.~8, a.~4}, he deepens that teaching by uniting generic spiritual conformity of soul with angel and the instrumental efficacy of Christ's deified humanity pouring light and joy into angelic minds. The angels do not receive redeeming forgiveness from Christ's cross, since faithful angels never sinned; yet through his sacred humanity they receive accidental joy, fuller illumination into the mystery of redemption, and governance under one Head.

\subsection{Angels in the Mysteries of Christ's Life and Exaltation}

From the conception of the Word to his return in glory, the Gospels frame every decisive mystery of Christ's life with angelic ministry. At the Nativity, an angel of the Lord stands by the shepherds in the brightness of God announcing ``good tidings of great joy that shall be to all the people,'' and suddenly ``a multitude of the heavenly army'' fills the night sky singing, ``Glory to God in the highest: and on earth peace to men of good will'' (Luke 2:8--15). As Pope St.~Leo the Great preached on the feast (\emph{Sermon~21 on the Nativity}, \S2), ``The mystery of the Incarnation is a fitting theme for joy both to angels and to men.'' Even Jesus himself, Pseudo-Dionysius marvels (\CH{IV.4}), honored the hierarchical order he established, submitting to the Father's dispositions communicated through angels in the flight into Egypt and the return to Israel (Matt 1:20; 2:13, 19).

In the wilderness, Satan turns Scripture against the Lord, citing the guardian-angel promise of Psalm 90:11--12 to lure Christ into presuming upon God; Christ refuses to tempt the Lord, commands Satan away, ``and behold angels came and ministered to him'' (Matt 4:1--11). They come not as superior guardians guiding a fallible wayfarer (\ST{I, q.~113, a.~4, ad~1}), but as loyal servants honoring their victorious King. On the night of his betrayal, the Lord who reminds Peter that his Father would instantly send ``more than twelve legions of angels'' (Matt 26:53) accepts instead the quiet ministry of ``an angel from heaven, strengthening him'' in his agony in Gethsemane (Luke 22:43).

At the empty tomb, an angel whose countenance is ``as lightning and his raiment as snow'' rolls back the stone, strikes the guards like dead men, and proclaims: ``He is not here. For he is risen, as he said'' (Matt 28:2--7; Luke 24:4--8). Inside the sepulchre, Mary Magdalene beholds ``two angels in white, sitting, one at the head, and one at the feet, where the body of Jesus had been laid'' (John 20:12). Forty days later, the Ascension carries human flesh above every created spirit. Pope St.~Leo the Great contemplates that triumph in his second Ascension homily (\emph{Sermon~74}, \S1):

\begin{quotation}
We commemorate and duly venerate that day on which the Nature of our humility in Christ was raised above all the host of heaven, over all the ranks of angels, beyond the height of all powers, to sit with God the Father.
\end{quotation}

St.~John Chrysostom draws out St.~Paul's words on that session at the Father's right hand (Eph 1:21): ``Whatever there be in Heaven, He has become above all'' (\emph{Homily~3 on Ephesians}). At the end of time, Christ will return ``in his majesty, and all the angels with him'' (Matt 25:31), descending ``with commandment and with the voice of an archangel and with the trumpet of God'' (1 Thess 4:15 [Vulg.\ 4:16]; 2 Thess 1:7) to send his angels as reapers separating the wicked from the just (Matt 13:39--49).

\subsection{The Annunciation to the Blessed Virgin (\ST{III, q.~30, aa.~1--4})}

No earthly encounter between angel and human being equals the embassy of Gabriel at Nazareth (Luke 1:26--38), treated in four articles of the \emph{Tertia Pars} (\ST{III, q.~30, aa.~1--4}).

First, the Incarnation was fittingly announced to Mary before its accomplishment (\ST{III, q.~30, a.~1}): that she might conceive Christ in her mind by faith before conceiving him in her flesh (St.~Augustine, \emph{De Sancta Virginitate}~III: ``Mary is more blessed in receiving the faith of Christ, than in conceiving the flesh of Christ''); that she might be a surer witness instructed from heaven; that she might offer free obedience; and to signify the spiritual marriage between the Son of God and mankind. ``Wherefore in the Annunciation the Virgin's consent was besought in lieu of that of the entire human nature.'' Divine predestination is fulfilled ``without the causality of our will; not without its consent'' (ad~1), and Mary's humility kept her from presuming such election for herself even though she already believed explicitly in the coming Incarnation (ad~2).

Second, an angel fittingly bore the message (\ST{III, q.~30, a.~2}), honoring the Dionysian law whereby heavenly mysteries reach men through angels (\CH{IV}), reversing Eve's dialogue with the serpent (St.~Bede), and matching the kinship of virginity with angelic life (Pseudo-Jerome: ``to live in the flesh and not according to the flesh is not an earthly but a heavenly life''). Superior to angels in election, Mary received angelic instruction as an earthly wayfarer (ad~1). Quoting St.~Maximus, Aquinas observes that the angels knew the Incarnation would occur, yet ``it was not given to them to trace the manner of our Lord's conception, nor how it was that He remained whole in the Father, whole throughout the universe, and was whole in the narrow abode of the Virgin'' (ad~3). Reconciling St.~Gregory's phrase ``one of the highest angels'' (\emph{Homily~34}, \S9) with Pseudo-Dionysius's rule that the first triad performs no external mission (\ST{I, q.~112, a.~2}), Aquinas concludes that Gabriel was not of the Seraphim, yet ``it is therefore sufficiently credible that he was the highest of the archangels'' (ad~4). His name matched his mission: ``Gabriel means `Power of God.' This message therefore was fittingly brought by the `Power of God,' because the Lord of hosts and mighty in battle was coming to overcome the powers of the air'' (Gregory, \emph{Homily~34}, \S9).

Third, Gabriel appeared in an external bodily vision (\ST{III, q.~30, a.~3}; \ST{I, q.~51, a.~2}), fitting the visible coming of the invisible God, refreshing Mary's senses with her mind, and affording greater certainty than Joseph's dream---as St.~John Chrysostom observes (\emph{Homily~4 on Matthew}), Gabriel ``came to the Virgin not in her sleep, but visibly'' (resp.\ and ad~1--2). Addressing Mary's disturbance (ad~3), Aquinas appeals to the discernment rule of St.~Athanasius's \emph{Life of Antony} (\S35, quoted in Section~\ref{sec:patristic-witness}): unlike demonic tumult, a holy angelic vision arrives with quietness and swiftly dispels human awe with joy and courage, just as Gabriel calmed Zachary and now reassures Mary (Luke 1:30). Her trouble sprang from virginal modesty and humble wonder ``at his saying'' (Luke 1:29).

Fourth, Gabriel spoke in becoming order (\ST{III, q.~30, a.~4}; Rom 13:1): winning attention with an unprecedented greeting---\emph{Full of grace} (worthiness), \emph{The Lord is with thee} (conception), \emph{Blessed art thou among women} (honor)---unfolding the Child's birth, kingship, and conception by the Holy Ghost, and leading her to consent through the sign of Elizabeth and divine omnipotence. Unlike Zachary's doubting ``Whereby shall I know this?'' (Luke 1:18), Mary's ``How shall this be done, because I know not man?'' (Luke 1:34) voiced humble wonder rather than unbelief (St.~Ambrose on Luke 1:34), issuing in her free surrender: ``Behold the handmaid of the Lord: be it done to me according to thy word'' (Luke 1:38).

\subsection{Mary, Queen of Angels}

In his fourth homily on \emph{Missus est} (\S8), St.~Bernard of Clairvaux arrests the moment between Gabriel's question and Mary's reply, picturing heaven and earth holding their breath before a young Jewish virgin:

\begin{quotation}
The angel awaits your reply, for it is time that he should return to God, Who sent him. We, too, are waiting, O Lady, for a word of mercy---we, who are groaning under the sentence of condemnation. See, the price of our salvation is offered to you; if you consent, we shall at once be delivered. By the Eternal Word of God we were all created, and behold we die. By your short answer we shall be refreshed and recalled to life.
\end{quotation}

From that \emph{fiat} flows the Church's confession of Mary as Queen of Angels. In the Dogmatic Constitution \emph{Lumen gentium} (\S\S53, 56, 59, 66), the Second Vatican Council teaches that Mary, redeemed in a more sublime manner in view of her Son's merits and endowed with the supreme dignity of Mother of God, far surpasses all other creatures both in heaven and on earth while remaining genuinely a daughter of Adam and a creature; taken up body and soul into heavenly glory, she is exalted by the Lord as Queen of the universe so as to be more fully conformed to her Son, Lord of lords. Pope Pius XII, in the encyclical \emph{Ad caeli Reginam} (1954), grounds this universal queenship over angels and men alike in her divine motherhood and in her singular partnership with Christ the King in the work of redemption.

At Bethlehem, the heavenly hosts themselves ratify that order of grace, gathering in awe around the Mother and her Child. St.~Alphonsus Liguori meditates on that contrast in \emph{The Incarnation, Birth, and Infancy of Jesus Christ}: ``The Son of God is born humble and poor in a stable. There indeed the angels of heaven acknowledge him, singing, Glory to God in the highest; but the inhabitants of the earth, for whose salvation Jesus was born, leave him neglected.'' Within the single City of God whose ranks are perfected not by created nature alone but by divine grace (\ST{I, q.~108, a.~8}), the redeemed Mother of the Word Incarnate stands above Seraphim and Cherubim alike---never a rival to her Son's unique mediation, yet first in grace and glory beneath him whom all the angels of God adore.
